\section{BUG-004: Hook Block Silent Failure}
\label{sec:bug-004}

\subsection*{Metadata}
\begin{tabular}{ll}
\textbf{Issue ID} & BUG-004 \\
\textbf{Severity} & Medium \\
\textbf{Category} & Bug Fix \\
\textbf{Affected File} & \srcfile{src/hooks.ts:120-145} \\
\textbf{Branch} & \branchref{fix/BUG-004-hook-block-silent-failure} \\
\textbf{Changes} & \branchdiff{fix/BUG-004-hook-block-silent-failure} \\
\end{tabular}

\subsection*{Executive Summary}
The hook execution system in \texttt{src/hooks.ts} is designed as a pluggable security layer: hooks can \texttt{allow}, \texttt{block}, or \texttt{modify} agent actions. The \texttt{runHooks} function at line 120 iterates over all registered hooks and returns the first result with \texttt{action === "block"} or \texttt{action === "modify"}. If a hook throws an error during execution (lines 130-141), the error is caught, logged to console via \texttt{console.error}, and \textbf{silently swallowed} — the iteration continues to the next hook.

This means that if a critical security hook (e.g., \texttt{pre\_tool\_use} that validates file access) fails due to a runtime error (missing script, permission denied, invalid JSON output, timeout), the agent \textbf{continues operating as if the hook passed}. The error is only visible in the console log, which the agent itself does not read programmatically. The security policy is silently bypassed.

The bug is at \texttt{src/hooks.ts:138-141}, where \texttt{catch} discards the error and allows the for-loop to continue, ultimately returning \texttt{\{ action: "allow" \}} if all hooks either pass or error.

---

\subsection*{Root Cause Analysis}
\#\#\# The Code



\begin{lstlisting}[language=TypeScript]
// hooks.ts:120-145
export async function runHooks(
    hooks: HookDefinition[] | undefined,
    agentDir: string,
    input: Record<string, any>,
): Promise<HookResult> {
    if (!hooks || hooks.length === 0) {
        return { action: "allow" };
    }

    for (const hook of hooks) {
        try {
            const result = await executeHook(hook, agentDir, input);
            if (result.action === "block") {
                return result;
            }
            if (result.action === "modify") {
                return result;
            }
        } catch (err: any) {
            console.error(`Hook error: ${err.message}`);    // ← Only logs to console
            // Hook errors don't block execution by default      ← Silently continues
        }
    }

    return { action: "allow" };  // ← If all hooks errored, returns allow!
}

\end{lstlisting}



\#\#\# The executeHook Function

The \texttt{executeHook} function (lines 44-117) can fail for many reasons:



\begin{lstlisting}[language=TypeScript]
// hooks.ts:44-117
async function executeHook(
    hook: HookDefinition,
    agentDir: string,
    input: Record<string, any>,
): Promise<HookResult> {
    // Programmatic hooks
    if (typeof hook._handler === "function") {
        try {
            const result = await hook._handler(input);
            return result ?? { action: "allow" };
        } catch (err: any) {
            throw new Error(`Programmatic hook failed: ${err.message}`);
        }
    }

    return new Promise((promiseResolve, reject) => {
        // Path traversal guard
        const resolvedScript = resolve(scriptPath);
        const allowedBase = resolve(baseDir);
        if (!resolvedScript.startsWith(allowedBase + "/") && resolvedScript !== allowedBase) {
            reject(new Error(`Hook "${hook.script}" escapes its base directory`));
            return;
        }

        const child = spawn("sh", [resolvedScript], { ... });
        // ... stdout/stderr, stdin, timeout ...
        const timeout = setTimeout(() => {
            child.kill("SIGTERM");
            reject(new Error(`Hook "${hook.script}" timed out after 10s`));
        }, 10_000);

        child.on("error", (err) => { reject(new Error(`...failed to start: ${err.message}`)); });
        child.on("close", (code) => {
            if (code !== 0) {
                reject(new Error(`...exited with code ${code}: ${stderr.trim()}`));
                return;
            }
            try {
                const result = JSON.parse(stdout.trim()) as HookResult;
                promiseResolve(result);
            } catch {
                // If hook doesn't return JSON, treat as allow
                promiseResolve({ action: "allow" });
            }
        });
    });
}

\end{lstlisting}



\#\#\# Failure Scenarios That Bypass Security

\textbf{Scenario A: Pre-tool-use security hook script is missing}



\begin{lstlisting}[language=TypeScript]
Hook definition: { script: "check-permissions.sh" }
check-permissions.sh was deleted or renamed

executeHook: spawn fails → child.on("error") → reject
runHooks catch: "Hook 'check-permissions.sh' failed to start: ENOENT"
→ Console log only → continues to next hook → returns allow
→ Agent executes tool without permission check

\end{lstlisting}



\textbf{Scenario B: Security hook times out}



\begin{lstlisting}[language=TypeScript]
Hook script hangs (network call, infinite loop, deadlock)

executeHook: setTimeout → child.kill("SIGTERM") → reject
runHooks catch: "Hook 'check-permissions.sh' timed out after 10s"
→ Console log only → continues → returns allow
→ Agent executes tool during the 10-second hang window AND after

\end{lstlisting}



\textbf{Scenario C: Security hook returns invalid JSON}



\begin{lstlisting}[language=TypeScript]
Hook script: echo "not valid json"

executeHook: JSON.parse throws → catches → returns { action: "allow" }
→ Security hook that should have blocked instead passes silently

\end{lstlisting}



Scenario C is actually in \texttt{executeHook} itself (line 110-114), where JSON parse failure defaults to \texttt{allow}. This compounds the problem.

\#\#\# Why This Is Medium Severity

1. \textbf{Security hooks are the last line of defense}: If a security-critical hook fails, the agent operates with no guardrails
2. \textbf{Silent failure}: Only visible in \texttt{console.error}, which is not programmatically checked
3. \textbf{Compound risk}: A single faulty hook script (or a filesystem glitch) disables the entire hook chain without warning
4. \textbf{False sense of security}: Deployments configure hooks believing they provide protection, but a runtime failure silently removes that protection

---

\subsection*{Fix Implementation}
\#\#\# Option A: Fail-Closed on Hook Error

Change \texttt{runHooks} to block execution when a hook errors:



\begin{lstlisting}[language=TypeScript]
// hooks.ts — fail-closed runHooks
export async function runHooks(
    hooks: HookDefinition[] | undefined,
    agentDir: string,
    input: Record<string, any>,
): Promise<HookResult> {
    if (!hooks || hooks.length === 0) {
        return { action: "allow" };
    }

    for (const hook of hooks) {
        try {
            const result = await executeHook(hook, agentDir, input);
            if (result.action === "block") {
                return result;
            }
            if (result.action === "modify") {
                return result;
            }
        } catch (err: any) {
            console.error(`Hook error: ${err.message}`);
            // FAIL CLOSED: hook errors block execution by default
            return {
                action: "block",
                reason: `Hook "${hook.description || hook.script}" failed: ${err.message}. Action blocked to maintain security policy.`,
            };
        }
    }

    return { action: "allow" };
}

\end{lstlisting}



\textbf{What this fixes:}
1. Any hook error results in \texttt{block} — fail-closed behavior
2. Security policy cannot be silently bypassed
3. Error message is propagated to the caller (not just console)

\#\#\# Option B: Configurable Failure Mode

Add an \texttt{onError} field to \texttt{HookDefinition}:



\begin{lstlisting}[language=TypeScript]
// hooks.ts — configurable failure mode
interface HookDefinition {
    script: string;
    description?: string;
    baseDir?: string;
    _handler?: (ctx: Record<string, any>) => Promise<HookResult> | HookResult;
    onError?: "allow" | "block";  // NEW: configure failure behavior
}

// In runHooks
try {
    const result = await executeHook(hook, agentDir, input);
    // ...
} catch (err: any) {
    console.error(`Hook error: ${err.message}`);
    if (hook.onError === "block") {
        return { action: "block", reason: `Hook failed: ${err.message}` };
    }
    // onError === "allow" (or undefined) — continue (current behavior)
}

\end{lstlisting}



This preserves backward compatibility while allowing security hooks to opt into fail-closed behavior.

---

\subsection*{Affected Files}
\begin{itemize}
    \item \srcfile{pocs/BUG-004-hook-block-silent-failure/README.md}
    \item \srcfile{pocs/BUG-004-hook-block-silent-failure/validation.ts}
    \item \srcfile{src/hooks.ts} — primary affected file
    \item \srcfile{src/index.ts}
    \item \srcfile{src/sdk.ts}
    \item \srcfile{test/hooks.test.ts}
\end{itemize}

\subsection*{Verification}
The fix is verified through unit tests and integration tests that confirm the corrected behavior:

\begin{lstlisting}[language=TypeScript]
// verification/bug-004-verify.ts
import { describe, it } from "node:test";
import assert from "node:assert";

interface HookResult {
    action: "allow" | "block" | "modify";
    reason?: string;
}

interface HookDefinition {
    script: string;
    description?: string;
    _handler?: () => Promise<HookResult>;
}

// Fixed implementation
async function executeHook(hook: HookDefinition): Promise<HookResult> {
    if (hook._handler) {
        return hook._handler();
    }
    // Simulate script-based hook
    throw new Error(`Hook "${hook.script}" failed to start: ENOENT`);
}

function runHooksFixed(hooks: HookDefinition[]): Promise<HookResult> {
    return (async () => {
        if (!hooks || hooks.length === 0) {
            return { action: "allow" };
        }

        for (const hook of hooks) {
            try {
                const result = await executeHook(hook);
                if (result.action === "block" || result.action === "modify") {
                    return result;
                }
            } catch (err: any) {
                console.error(`Hook error: ${err.message}`);
                // FAIL CLOSED
                return {
                    action: "block",
                    reason: `Hook "${hook.description || hook.script}" failed: ${err.message}`,
                } as HookResult;
            }
        }

        return { action: "allow" };
    })();
}

// Old (buggy) implementation for comparison
function runHooksOld(hooks: HookDefinition[]): Promise<HookResult> {
    return (async () => {
        if (!hooks || hooks.length === 0) {
            return { action: "allow" };
        }

        for (const hook of hooks) {
            try {
                const result = await executeHook(hook);
                if (result.action === "block" || result.action === "modify") {
                    return result;
                }
            } catch (err: any) {
                console.error(`Hook error: ${err.message}`);
                // Hook errors don't block execution by default
            }
        }

        return { action: "allow" };
    })();
}

describe("BUG-004: Hook Block Silent Failure", () => {
    it("OLD BEHAVIOR: should silently allow when hook fails", async () => {
        const hooks: HookDefinition[] = [
            { script: "security-check.sh", description: "Security check" },
        ];

        const result = await runHooksOld(hooks);
        assert.strictEqual(result.action, "allow",
            "OLD: Failing hook should return allow (BUG)");
    });

    it("FIXED: should block when hook fails", async () => {
        const hooks: HookDefinition[] = [
            { script: "security-check.sh", description: "Security check" },
        ];

        const result = await runHooksFixed(hooks);
        assert.strictEqual(result.action, "block",
            "FIXED: Failing hook should return block");
        assert.ok(result.reason?.includes("failed"),
            "FIXED: Reason should mention the failure");
    });

    it("should still allow when hooks pass", async () => {
        const hooks: HookDefinition[] = [
            {
                script: "passing-hook.sh",
                description: "Passing check",
                _handler: async () => ({ action: "allow" } as HookResult),
            },
        ];

        const result = await runHooksFixed(hooks);
        assert.strictEqual(result.action, "allow",
            "Passing hook should return allow");
    });

    it("should still block when hook explicitly blocks", async () => {
        const hooks: HookDefinition[] = [
            {
                script: "blocking-hook.sh",
                description: "Blocking check",
                _handler: async () => ({ action: "block", reason: "Not allowed" } as HookResult),
            },
        ];

        const result = await runHooksFixed(hooks);
        assert.strictEqual(result.action, "block",
            "Explicit block should still work");
    });

    it("should handle mixture of passing and failing hooks", async () => {
        const hooks: HookDefinition[] = [
            { script: "failing-hook.sh", description: "Failing check" },
            {
                script: "passing-hook.sh",
                description: "Passing check",
                _handler: async () => ({ action: "allow" } as HookResult),
            },
        ];

        const result = await runHooksFixed(hooks);
        assert.strictEqual(result.action, "block",
            "First hook fails → should block before reaching passing hook");
    });

    it("should return allow for empty hooks", async () => {
        const result = await runHooksFixed([]);
        assert.strictEqual(result.action, "allow",
            "Empty hooks should return allow");
    });
});

\end{lstlisting}

